\documentclass[../Main.tex]{subfiles}


% Define styles
\tikzstyle{process} = [rectangle, rounded corners, minimum width=3cm, minimum height=1cm, text centered, draw=black, fill=white]
\tikzstyle{arrow} = [thick,->,>=stealth]
\tikzstyle{data} = [cylinder, shape border rotate=90, minimum height=1.5cm, minimum width=1cm, text centered, draw=black, fill=white]
\tikzstyle{nodecircle} = [circle, draw=black, minimum size=1cm, text centered]
    
\begin{document}

\begin{figure}
    \centering
    \resizebox{\textwidth}{!}{
\begin{tikzpicture}[node distance=2cm]

% Section 3: Travel Survey
\node (survey) [data] {Microdata};


% Section 4: Structural Learning
\node (learning) [process, right=1cm of survey] {Structural learning};

% Structural learning graph (Bayesian network)

\begin{scope}[scale=0.48, every node/.style={transform shape}]
\node (A) [nodecircle, below=1.5cm of learning, xshift=-1.5cm, yshift = 8cm] {A};
\node (B) [nodecircle, above right=1.5cm of A] {B};
\node (C) [nodecircle, below right=1.5cm of B] {C};
\node (D) [nodecircle, below=1.5cm of B] {D};
\node (E) [nodecircle, below=1.5cm of C] {E};
\draw [arrow] (A) -- (B);
\draw [arrow] (B) -- (C);
\draw [arrow] (B) -- (D);
\draw [arrow] (C) -- (E);
\draw [arrow] (D) -- (E);

\node (structural_group) [draw=black, dashed, fit={(learning) (A) (B) (C) (D) (E)}, inner sep=0.5cm] {};
\end{scope}
% Posterior probability distributions (bar chart)
\node (posterior) [process, right=2cm of learning] {Posterior probability  distributions};
% Bar chart
\begin{scope}[scale=0.8, every node/.style={transform shape}]
\node (chart) [draw=none, below=-5cm of posterior] {
    \begin{tikzpicture}
        \draw[thick] (0,0) -- (3,0); % x-axis
        \draw[thick] (0,0) -- (0,2.5); % y-axis
        \foreach \x in {0.5, 1, 1.5, 2} {
        \pgfmathsetmacro{\height}{abs(rand*2)} % Ensure positive height
            \draw[fill=blue!50] (\x-0.2,0) rectangle (\x+0.2, {\height});
        }
        \node at (1.5, -0.5) {Categories}; % x-axis label
        \node[rotate=90] at (-0.5, 1.25) {Probability}; % y-axis label
    \end{tikzpicture}
};
\node (posterior_group) [draw=black, dashed, fit={(posterior) (chart)}, inner sep=0.5cm] {};
\end{scope}
% Sample activity and trip chains
%\node (sample) [process, above=1.5cm of posterior] {Sample activity and \\ trip chains};
%\draw [arrow] (posterior.north) -- (sample.south);
%
% Section 5: Evaluate predicted samples
\node (predict) [process, right=2cm of posterior] {Sample generative agents};
\node (evaluate) [data, right=2cm of predict] {Expanded Data};
%\draw [arrow] (learning.south) -- (evaluate.north);

% Section Labels
%\node at ($(survey.west) + (-1,0)$) [align=center] {\textbf{Section 3}};
%\node at ($(learning.north) + (0,1.5)$) [align=center] {\textbf{Section 4}};
%\node at ($(evaluate.south) + (0,-1)$) [align=center] {\textbf{Section 5}};

% Connect test set to learning
\draw [arrow] (survey.east) -- ++(0.5,0) |- (learning.west);
\draw [arrow] (learning.east) -- ++(0.5,0) |- (posterior.west);
\draw [arrow] (posterior.east) -- ++(0.5,0) |- (predict.west);
\draw [arrow] (predict.east) -- ++(0.5,0) |- (evaluate.west);
\end{tikzpicture}}
\caption{Flow Process of the Bayesian Network}
    \label{fig:flow}
    
    \end{figure}
\end{document}
